\textit{Q1--Q10 form a 100-point multiple-choice section. The supplied paper
does not state individual point allocations.}

For this multiple-choice problem, choose only one correct answer. Write the
chosen answer and the original problem number on the answer sheet.

A small satellite (a luminous sphere) that orbits the Earth, in an equatorial
circular orbit at the altitude $h$, is photographed with a camera on the
ground, the lens of which is a converging lens with focal length $f$.
The photography is made when the satellite aligns with the camera and the
centre of the Earth.

It is given: the mass of the Earth, $M$; the gravitational constant, $G$;
the radius of the Earth, $R$. During the photographic exposure the proper
rotation of the Earth is neglected.

The length of the image on the photograph, after developing the photographic
plate, placed in the focal plane of the objective lens, as indicated by the
drawing, if the exposure time of the camera was $\Delta t$, is given by the
expression
\ioaafig{2022_TH_Q04-f1.pdf}
\begin{description}
\item[(a)] $l=\dfrac{f}{h-f}\sqrt{G\dfrac{M}{R+h}}\,\Delta t$.
\item[(b)] $l=\dfrac{f}{h+f}\sqrt{G\dfrac{M}{R+h}}\,\Delta t$.
\item[(c)] $l=\dfrac{f}{h-f}\sqrt{G\dfrac{M}{h-R}}\,\Delta t$.
\item[(d)] $l=\dfrac{f}{h-f}\sqrt{2G\dfrac{M}{R+h}}\,\Delta t$.
\end{description}
